\chapter{相关研究}
\label{chap:related}

\section{可读性公式}

Flesch 阅读难度公式\cite{flesch1948}用平均句长和平均音节数估计英文文本的难度，
至今仍被广泛使用。中文方面，研究者以字频、词频和句长为特征，
提出了多种面向教材与新闻的可读性公式\cite{wang2020}。
这些方法的共同点是只看文字本身。

\section{排版与阅读}

Tinker 与 Paterson 在 20 世纪 20 年代系统研究了行长、字号与行距对阅读速度的影响\cite{tinker1929}；
近年的研究把同样的问题搬到了屏幕上\cite{dyson2004}。
中文排版的规范性要求见于《中文排版需求》\cite{w3c2020clreq}，
但它给出的是排版规则，并未量化这些规则对读者的影响。

\section{小结}

已有工作要么只看文字，要么只看版面，
缺少一个把两者放在同一框架下评估的方法。这正是本文要补上的部分。
